\documentclass[a4paper,12pt]{article}
\usepackage[portuges]{babel}
\usepackage[latin1]{inputenc}
\usepackage{graphicx}
\usepackage{amsmath,amsfonts,mathrsfs,amssymb}
\usepackage{multicol}
\usepackage{fancyhdr}
\usepackage{makeidx}
\numberwithin{equation}{section}
\usepackage{geometry}
\geometry{verbose,tmargin=2cm,bmargin=1cm,lmargin=2cm,rmargin=2cm}
\usepackage[pdftex,plainpages=false,pdfpagelabels,pagebackref,colorlinks=true,
citecolor=black,linkcolor=black,urlcolor=black,filecolor=black,bookmarksopen=true]{
hyperref}
\usepackage{sectsty}
\sectionfont{\large}
\subsectionfont{\normalsize}


\setlength{\textheight}{240mm}
\setlength{\headheight}{28pt}
\renewcommand{\headrulewidth}{0.5pt}
\renewcommand{\footrulewidth}{0.5pt}
\renewcommand{\baselinestretch}{1.2}


%Atalhos para nota��o de conjuntos
\newcommand{\cqd}{{{\hfill $\rule{2.5mm}{2.5mm}$}}\vspace{1cm}\\}
\newcommand{\sen}{\mathrm{sen}}
\newcommand{\ds}{\displaystyle}
\newcommand{\R}{\mathbb{R}}
\newcommand{\C}{\mathbb{C}}
\newcommand{\N}{\mathbb{N}}
\newcommand{\Z}{\mathbb{Z}}
\newcommand{\Q}{\mathbb{Q}}
\newcommand{\I}{\mathbb{I}}


%Comandos Atalhos para opera��es matem�ticas
\newcommand{\deli}[2]{\ensuremath\frac{\partial#1}{\partial#2}}
\newcommand{\deri}[2]{\ensuremath\frac{d#1}{d#2}}
\newcommand{\nt}[4]{\ensuremath{ \int_{#1}^{#2} #3 \, \mbox{d}#4}}

%Atalhos para nota��es matem�ticas, fun��es, conjuntos, etc.
\newcommand{\raiotres}{\sqrt{x^2 + y^2 + z^2}}
\newcommand{\coordn}[1]{(#1_1,\dotsc,#1_n)}
\newcommand{\normapi}[2]{\langle #1,#2 \rangle}
\newcommand{\norma}[1]{||#1||}
\newcommand{\prodint}[2]{\langle #1,#2 \rangle}
\newcommand{\setn}[1]{(#1_1,\dotsc,#1_n)}
\newcommand{\traco}{{\text{Tr}}}
\newcommand{\posto}{{\text{posto}}}
\newcommand{\Ima}{{\text{Im}}}

%Atalhos para ambientes matem�ticos
\newcommand{\exercicio}[1]{\vspace{0.2cm}\noindent\textbf{Exerc�cio #1. }\ }
\newcommand{\lema}[1]{\vspace{0.2cm}\noindent\textbf{Lema #1. }\ }
\newcommand{\propriedade}[1]{\vspace{0.2cm}\noindent\textbf{Propriedade #1:}\ }
\newcommand{\definicao}{\vspace{0.2cm}\noindent\textbf{Defini��o:}\ }
\newcommand{\solucao}{\vspace{0.2cm}\noindent\textbf{Solu��o:}\ }
\newcommand{\exemplo}[1]{\vspace{0.2cm}\noindent\textbf{Exemplo #1 :}\ }
\newcommand{\demonstracao}{\vspace{0.2cm}\noindent\textbf{Demonstra��o:}\ }
\newcommand{\prova}{\vspace{0.2cm}\noindent\textbf{Prova:}\ }

\pagestyle{fancy}
\fancyhead[LO,L]{\textit \normalsize \nouppercase \leftmark}
\fancyhead[RO,R]{\textit \normalsize \nouppercase \rightmark}
\fancyfoot[LO,L]{\itshape{Osvaldo Pereira}}
\fancyfoot[RO,R]{Polin�mios de Legendre}
\begin{document}
\tableofcontents
\section{Polin�mios de Legendre}


\subsection{C�lculo de (2n)!}
Seja:
%
%
\[f_n(t) = \frac{1}{2^nn!}\,\frac{d^n}{dt^n}(t^2 - 1)^n\]
%
%
Expandindo a fun��o $u(t) = t^2 - 1$ em termos do teorema binomial, temos que:
%
%
\[u^n = (t^2 - 1)^n = \sum_{k=0}^{n}\binom{n}{k}t^{2n - 2k}(-1)^{k}\]
%
%
Derivando uma vez encontramos que:
%
%
\[(u^n)' = \sum_{k=0}^{n-1}\binom{n}{k}(2n - 2k)t^{2n - 2k - 1}(-1)^{k} \]
%
%
Derivando pela segunda vez:
%
%
\[(u^n)'' = \sum_{k=0}^{n-2}\binom{n}{k}(2n - 2k)(2n - 2k - 1)t^{2n - 2k - 2}(-1)^{k} \]
%
%
Derivando m vezes, temos que:
%
%
\[(u^n)^{(m)} = \sum_{k=0}^{n}\binom{n}{k}(2n - 2k)(2n - 2k - 1)...(2n-2k- m + 1)t^{2n - 2k - m}(-1)^{k} \]
%
%
Derivando n vezes:
%
%
\[(u^n)^{(n)} = \sum_{k=0}^{N}\binom{n}{k}(2n - 2k)(2n - 2k - 1)...(n - 2k)t^{n - 2k}(-1)^{k} \]
%
%
Sendo $N = n/2$ se $n$ for par, e $N = (n - 1)/2$ se $n$ for �mpar. Derivando $2n$ vezes, a somat�ria tem apenas um termo, sendo o �nico valor poss�vel: $k = 0$
%
%
\[(u^n)^{(2n)} = \sum_{k = 0}^{0}\binom{n}{k}(2n - 2k)(2n - 2k - 1)...(1 - 2k)t^{- 2k}(-1)^{k} \]
%
%
Ou seja:
%
%
\[(u^n)^{(2n)} = \binom{n}{k}(2n - 0)(2n - 0 - 1)\ldots(1 - 0)t^{0}(-1)^{0} = (2n)(2n - 1)\ldots1 = (2n)! \]



\newpage
\subsection{Normaliza��o}

Seja:
%
%
\[f_n(t) = \frac{1}{2^nn!}\,\frac{d^n}{dt^n}(t^2 - 1)^n\]
%
%
Fazendo $u = t^2 - 1$, temos que:
%
%
\[p_n(t) = 2^nn!f_n(t) = \frac{d^n}{dt^n}u^n = (u^n)^{(n)}\]
%
%
Calculando:
%
%
\[ ||p_n||^2 = \int_{-1}^{1} p_n^2(t)\, dt = \int_{-1}^{1} (u^n)^{(n)}(u^n)^{(n)}\, dt \]
%
%
Notemos que $u(t)$ se anula no intervalo $[-1,1]$, assim como todas suas n-�simas derivadas. De forma que, integrando uma vez por partes, encontramos que:
%
%
\[I_{n,0} = \int_{-1}^{1} (u^n)^{(n)}(u^n)^{(n)}\, dt = (u^n)^{(n-1)}(u^n)^{(n)}\big|_{-1}^{1} - \int_{-1}^{1} (u^n)^{(n-1)}(u^n)^{(n+1)}\,dt\]
%
%
\[I_{n,1} = - \int_{-1}^{1} (u^n)^{(n-1)}(u^n)^{(n+1)}\,dt\]
%
%
Integrando novamente por partes, encontramos que:
%
%
\[I_{n,1} = - \int_{-1}^{1} (u^n)^{(n-1)}(u^n)^{(n+1)}\, dt = - (u^n)^{(n-2)}(u^n)^{(n+1)}\big|_{-1}^{1} + \int_{-1}^{1} (u^n)^{(n-2)}(u^n)^{(n+2)}\, dt\]
%
%
Integrando novamente:
%
%
\[I_{n,2} = \int_{-1}^{1} (u^n)^{(n-2)}(u^n)^{(n+2)}\, dt\]
%
%
Integrando m vezes por partes, encontramos que:
%
%
\[I_{n,m} = (-1)^m\int_{-1}^{1} (u^n)^{(n-m)}(u^n)^{(n+m)}\, dt\]
%
%
Integrando n vezes:
%
%
\[I_{n,n} = (-1)^n\int_{-1}^{1} (u^n)(u^n)^{(2n)}\, dt = (2n)!(-1)^n\int_{-1}^{1} (u^n)\, dt\]
%
%
Voltando a nota��o antiga:
%
%
\[I_{n,n} = (-1)^n\int_{-1}^{1} (u^n)(u^n)^{(2n)}\, dt = (2n)!\int_{-1}^{1} (1 - t^2)^n\, dt\]
%
%
Agora o problema se resume em calcular a seguinte integral:
%
%
\[I_n = \int_{-1}^{1} (1 - t^2)^n\, dt\]




\subsection{C�lculo da Integral}

Calculando a integral por partes n vezes, e depois simplificando a rela��o, encontra-se uma rela��o de recorr�ncia entre as integrais $In$ e $I_{n-1}$:

\begin{eqnarray}
\nonumber I_n &=& \int_{-1}^{1} (1 - t^2)^n\, dt \\
\nonumber     &=& \int_{-1}^{1} (1 - t^2)^{n-1} (1 - t^2)\, dt\\
\nonumber     &=& I_{n-1} - \int_{-1}^{1} (1 - t^2)^{n-1}t^2\, dt\\
\nonumber     &=& I_{n-1} + \frac{1}{2}\int_{-1}^{1} t(1 - t^2)^{n-1}\, d(1 - t^2)\\
\nonumber     &=& I_{n-1} + \frac{1}{2n}\int_{-1}^{1} t\, d[(1 - t^2)^n]\\
\nonumber     &=& I_{n-1} + \frac{1}{2n}\left(t(1 - t^2)^n\big|_{-1}^{1} - \int_{-1}^{1}(1 - t^2)^n\,dt\right)\\
\nonumber     &=& I_{n-1} - \frac{I_n}{2n}\\
\nonumber     &=& \frac{2n}{2n + 1}I_{n-1}
\end{eqnarray}
%
%
Abrindo a rela��o de recorr�ncia, encontra-se a seguinte rela��o:
%
%
\[I_n = \frac{2n}{2n + 1}\frac{2n - 2}{2n - 1}\frac{2n - 4}{2n - 3}\ldots \frac{2n - 2k}{2n - 2k + 1}I_{n-k}\]
%
%
Fazendo k = n, notando que $(2n)\ldots 2 = 2^n n!$ e que completando o denominador com $2^n n!$ encontramos $(2n)!$, temos ent�o que:
%
%
\begin{eqnarray}
\nonumber I_n &=& \frac{2n}{2n + 1}\frac{2n - 2}{2n - 1}\frac{2n - 4}{2n - 3}\ldots \frac{2}{3}I_{0}\\
\nonumber     &=& \frac{2^nn!}{2n + 1}\frac{1}{2n - 1}\ldots\frac{2}{3} I_{0}\\
\nonumber     &=& \frac{2^nn!}{2n + 1}\frac{2n}{(2n)(2n - 1)}\ldots\frac{2}{2.3} I_{0}\\
\nonumber     &=& \frac{2^nn!}{2n + 1}\frac{2^{n+1}n!}{(2n)!} I_{0}\\
\nonumber     &=& \frac{2^{2n+1}}{2n + 1}\frac{(n!)^2}{(2n)!} I_{0}
\end{eqnarray}
%
%
De forma que (lembrando que $I_0 = 1$):
%
%
\[I_{n,n} = (2n)!\int_{-1}^{1} (1 - t^2)^n\, dt = \frac{2^{2n+1}}{2n + 1}(n!)^2\]


\subsection{Motiva��o: normaliza��o}

A partir da rela��o encontrada no c�lculo da norma de $f_n$:
%
%
\[ ||f_n||^2 = \frac{1}{2^{2n}(n!)^2}\int_{-1}^{1} f_n^2(t)\, dt = \frac{1}{2^{2n}(n!)^2}\int_{-1}^{1} (u^n)^{(n)}(u^n)^{(n)}\, dt \]
%
%
De forma que:
%
%
\begin{eqnarray}
\nonumber ||f_n||^2 &=& \frac{1}{2^{2n}(n!)^2} \int_{-1}^{1} f_n^2(t)\, dt \\
\nonumber           &=& \frac{1}{2^{2n}(n!)^2}\int_{-1}^{1} (u^n)^{(n)}(u^n)^{(n)}\, dt\\
\nonumber           &=& \frac{1}{2^{2n}(n!)^2}\frac{2^{2n+1}}{2n + 1}(n!)^2\\
\nonumber           &=& \frac{2}{2n + 1}
\end{eqnarray}
%
%
Finalmente encontramos a seguinte rela��o:
%
%
\[||f_n|| = \sqrt{\frac{2}{2n + 1}}\]
%
%
E assim fica justificada a motiva��o para a escolha das seguintes fun��es (fun��es normalizadas):
%
%
\[e_n = \sqrt{\frac{2n + 1}{2}}f_n \to ||e_n|| = 1\]



\end{document}